
\section{SRForGAN experimental results}

The experimental setup is based on the same setup of the previous ForGAN section \ref{sec:bs experiment}.
The asset price trajectories are simulated over a fixed time horizon of one month, discretized into $N=22$ daily time steps.
For each experiment, $J=100{,}000$ independent paths are generated, ensuring sufficient sample diversity and statistical robustness.
The initial asset price $X_0$ is set for each trajectory to $0.0$, while the drift parameter $\mu$ is fixed at zero to isolate the model's ability to learn volatility-driven dynamics.
The volatility parameter $\sigma$ is sampled from a wide range $(0.1, 1.2)$ to capture varying market conditions, from near-deterministic to highly stochastic regimes.
The PDF to be learned refers to the $1$-th step ahead distribution, that is the distribution of asset prices at time $T + 1$ (or $N dt + 1$) given information up to time $T$.\\

Also in this case, the target distribution of the next value $c_{T + \tau}$  by using montecarlo simulations. 
The results are provided for the adversarial free SRForGAN architecture and also for the two 
hybrid approaches using the ForGAN and ForWGAN.\\



\begin{table}[H]
\centering
\caption{Generator Architecture and Hyperparameters Summary}
\label{tab:srforgan_architecture}
\begin{tabular}{lll}
\hline
\textbf{Parameter}         & \textbf{Generator}                  & \textbf{Critic (Discriminator)} \\ \hline
Input Dimension            & 32 (Latent) + 22 (Cond.)           & --                              \\
Hidden Layers              & [128, 256, 128, 64]                & --                              \\
Output Dimension           & 1                                  & --                              \\
Activation Function        & Leaky ReLU                         & --                              \\
Normalization              & BatchNormalization                 & --                              \\
Dropout                    & 0.0                                & --                              \\
\multicolumn{3}{c}{\textbf{Global Training Settings}}                                      \\ \hline
Batch Size                 & \multicolumn{2}{l}{32}                                            \\
max\_epochs                & \multicolumn{2}{l}{100}                                          \\
\hline
\end{tabular}
\end{table}


\begin{table}[h]
\centering
\begin{tabular}{l c}
\hline
\textbf{Metric} & \textbf{Value} \\
\hline
\label{tab:srforgan_results}
KS Test Acceptance (\%) & 100.00 \\
JS Distance & 0.1365 \\
Hellinger Distance & 0.1178 \\
TV Distance & 0.1099 \\
EMD Distance & 0.0018 \\
\hline
\end{tabular}
\caption{Evaluation metrics rounded to four decimal places}
\end{table}

\begin{figure}[H]
    \centering
    \includegraphics[width=1.0\textwidth]{SRForGAN/images/SRForGAN_samples.png}
    \caption{Samples obtained using fixed $X_0$, $\mu$ and $\sigma = [0.3, 0.5, 0.7, 0.9]$.}
    \label{srforgan_sample}
\end{figure}

It can be noticed that the SRForGAN architecture is able to learn the target distribution with a high degree of accuracy, as evidenced by the evaluation metrics and visual inspection of the generated samples. 
The KS test acceptance rate of 100\% indicates that the generated samples are statistically indistinguishable from the true distribution, while the low values of JS, Hellinger, TV, and EMD distances 
further confirm the model's effectiveness in capturing the underlying dynamics of the asset price trajectories. These results demonstrate
the potential of scoring rule-based generative models with respect to traditional adversarial approaches.\\
Because of these considerations further experiments have been conducted using more complicated distributions,
such as the Heston model, to assess the performance of the SRForGAN architecture in capturing complex stochastic processes.

\begin{comment}

%----------------------------------------------------------
%----------------------------------------------------------
By using a combined loss function made by both the energy score loss and the adversarial BCE or Wasserstein loss
the experiment has been repeated with the following architecture:


\begin{table}[H]
\centering
\caption{ForGAN Architecture and Hyperparameters Summary}
\label{tab:forgan_hybrid_architecture}
\begin{tabular}{lll}
\hline
\textbf{Parameter}         & \textbf{Generator}                  & \textbf{Critic (Discriminator)} \\ \hline
Input Dimension            & 32 (Latent) + 22 (Cond.)           & 1 (Data) + 22 (Cond.)           \\
Hidden Layers              & [128, 256, 128, 64]                & [128, 256, 256, 128, 64]        \\
Output Dimension           & 1                                  & 1                               \\
Activation Function        & Leaky ReLU                         & Leaky ReLU                      \\
Normalization              & BatchNormalization                 & None                            \\
Dropout                    & 0.0                                & 0.0                             \\
\multicolumn{3}{c}{\textbf{Global Training Settings}}                                      \\ \hline
Batch Size                 & \multicolumn{2}{l}{32}                                            \\
max\_epochs                & \multicolumn{2}{l}{100}                                          \\
\hline
\end{tabular}
\end{table}




\begin{table}[h]
\centering
\begin{tabular}{l c}
\hline
\textbf{Metric} & \textbf{Value} \\
\hline
\label{tab:hybrid_forgan_results}
KS Test Acceptance (\%) & 0.8530 \\
JS Distance & 0.2822 \\
Hellinger Distance & 0.2524 \\
TV Distance & 0.24029 \\
EMD Distance & 0.0102\\
\hline
\end{tabular}
\caption{Evaluation metrics rounded to four decimal places}
\end{table}
    


\begin{figure}[H]
    \centering
    \includegraphics[width=1.0\textwidth]{SRForGAN/images/Hybrid_ForGAN_samples.png}
    \caption{Samples obtained using fixed $X_0$, $\mu$ and $\sigma = [0.3, 0.5, 0.7, 0.9]$.}
    \label{hybrid_forgan_sample}
\end{figure}




\begin{table}[H]
\centering
\caption{ForWGAN Architecture and Hyperparameters Summary}
\label{tab:forwgan_hybrid_architecture}
\begin{tabular}{lll}
\hline
\textbf{Parameter}         & \textbf{Generator}                  & \textbf{Critic (Discriminator)} \\ \hline
Input Dimension            & 32 (Latent) + 22 (Cond.)           & 1 (Data) + 22 (Cond.)           \\
Hidden Layers              & [128, 256, 128, 64]                & [128, 256, 256, 128, 64]        \\
Output Dimension           & 1                                  & 1                               \\
Activation Function        & Leaky ReLU                         & Leaky ReLU                      \\
Normalization              & None                               & LayerNormalization              \\
Dropout                    & 0.0                                & 0.0                             \\
\multicolumn{3}{c}{\textbf{Global Training Settings}}                                      \\ \hline
Batch Size                 & \multicolumn{2}{l}{32}                                            \\
max\_epochs                & \multicolumn{2}{l}{100}                                          \\
\hline
\end{tabular}
\end{table}


\begin{table}[h]
\centering
\begin{tabular}{l c}
\hline
\textbf{Metric} & \textbf{Value} \\
\hline
\label{tab:hybrid_forwgan_results}
KS Test Acceptance (\%) & 0.9520 \\
JS Distance & 0.2653 \\
Hellinger Distance & 0.2348 \\
TV Distance & 0.2291 \\
EMD Distance & 0.0110\\
\hline
\end{tabular}
\caption{Evaluation metrics rounded to four decimal places}
\end{table}
   



\begin{figure}[H]
    \centering
    \includegraphics[width=1.0\textwidth]{SRForGAN/images/Hybrid_ForWGAN_samples.png}
    \caption{Samples obtained using fixed $X_0$, $\mu$ and $\sigma = [0.3, 0.5, 0.7, 0.9]$.}
    \label{hybrid_forwgan_sample}
\end{figure}

\end{comment}